% BHU PYQ -- Professional Exam Template
\documentclass[11pt,a4paper]{article}

\usepackage[a4paper,top=2.5cm,bottom=2.5cm,left=2.8cm,right=2.8cm,headheight=14pt]{geometry}
\usepackage{amsmath,amssymb}
\usepackage[T1]{fontenc}
\usepackage[utf8]{inputenc}
\usepackage{lmodern}
\usepackage{microtype}
\usepackage{fancyhdr}
\usepackage{enumitem}
\usepackage{listings}
\usepackage{hyperref}
\usepackage{lastpage}
\usepackage{parskip}

\hypersetup{colorlinks=true,urlcolor=black,linkcolor=black,
  pdftitle={CS-102: Object Oriented Programming -- BHU PYQ}}

\pagestyle{fancy}
\fancyhf{}
\renewcommand{\headrulewidth}{0.4pt}
\renewcommand{\footrulewidth}{0.4pt}
\fancyhead[L]{\small   \textit{Banaras Hindu University}}
\fancyhead[R]{\small   \textit{CS-102: Object Oriented Programming}}
\fancyfoot[C]{\small Page  \thepage\ of \pageref{LastPage}}

\lstset{
  basicstyle=\small tfamily,
  frame=single,
  breaklines=true,
  columns=flexible,
  keepspaces=true
}


\newlist{parts}{enumerate}{2}
\setlist[parts,1]{label=(\alph*),leftmargin=*,itemsep=4pt,topsep=3pt}
\setlist[parts,2]{label=(\roman*),leftmargin=*,itemsep=2pt,topsep=2pt}

\newcommand{\pts}[1]{\hfill   \textit{[#1]}}


\begin{document}

\begin{center}
  {\large\bfseries BANARAS HINDU UNIVERSITY}\\[2pt]
  {\normalsize B.Sc. (Hons.) Semester II Examination, 2023-24}\\[6pt]
  \rule{0.8\linewidth}{0.5pt}\\[6pt]
  {\large\bfseries Paper: CS-102 --- Object Oriented Programming}\\[4pt]
  {\normalsize Time: Three hours \hfill Maximum Marks: 70}
\end{center}

\rule{\linewidth}{0.4pt}

\smallskip



\noindent   \textit{Instructions: Attempt five questions including Question No. 1, which is compulsory. All undefined symbols have their usual meanings.}

\bigskip

\medskip


\noindent   \textbf{Question 1.} Attempt all questions: \pts{$2 imes9=18$}

\begin{parts}
  \item Differentiate between procedural programming and object-oriented programming.
  \item What is dynamic binding in C++?
  \item Explain the purpose of a destructor.
  \item What is a reference variable? Show its usage with an example.
  \item What is a copy constructor? Write its prototype.
  \item Can a virtual function be declared inline? Why or why not?
  \item Find the output of the following code segment:

\begin{lstlisting}[language=C++]
#include <iostream>
using namespace std;
class Test {
    static int count;
public:
    Test() { count++; }
    static int getCount() { return count; }
};
int Test::count = 0;
int main() {
    Test t1, t2, t3;
    cout << "Count = " << Test::getCount() << endl;
    return 0;
}
\end{lstlisting}
  \item Identify and explain the compile-time error in the following code block:

\begin{lstlisting}[language=C++]
#include <iostream>
using namespace std;
class Box {
    double width;
public:
    Box(double w) : width(w) {}
};
void printWidth(Box b) {
    cout << b.width << endl;
}
\end{lstlisting}
  \item What is the use of   \texttt{this} pointer in C++?
\end{parts}

\medskip


\noindent   \textbf{Question 2.}

\begin{parts}
  \item What is encapsulation? How is it implemented in C++? Compare private, protected, and public access specifiers. \pts{6}
  \item Explain the concept of inline functions. What are the rules and situations where inline function expansion is ignored by the compiler? \pts{4}
  \item Discuss function overloading with a suitable code example. \pts{4}
\end{parts}

\medskip


\noindent   \textbf{Question 3.}

\begin{parts}
  \item Explain the dynamic memory allocation operators   \texttt{new} and   \texttt{delete} in C++ with syntax and examples. Differentiate between   \texttt{delete} and   \texttt{delete[]}. \pts{6}
  \item What are friend functions and friend classes? Discuss their advantages and disadvantages with a suitable example. \pts{8}
\end{parts}

\medskip


\noindent   \textbf{Question 4.}

\begin{parts}
  \item What is constructor overloading? Write a C++ class representing a   \texttt{Complex} number with default, parameterized, and copy constructors. \pts{8}
  \item Write the rules for overloading operators in C++. How does overloading a binary operator using a friend function differ from overloading it using a member function? \pts{6}
\end{parts}

\medskip


\noindent   \textbf{Question 5.}

\begin{parts}
  \item Write a C++ program to overload the binary addition operator   \texttt{+} to add two   \texttt{Matrix} objects. \pts{8}
  \item Explain unary operator overloading with an example of overloading pre-increment   \texttt{++} and post-increment   \texttt{++} operators. \pts{6}
\end{parts}

\medskip


\noindent   \textbf{Question 6.}

\begin{parts}
  \item What are the different types of inheritance supported in C++? Explain the concept of virtual base class in multiple inheritance with a diagram and code snippet to show how ambiguity is resolved. \pts{8}
  \item Discuss the difference between compile-time polymorphism and runtime polymorphism. How is runtime polymorphism achieved in C++? \pts{6}
\end{parts}

\medskip


\noindent   \textbf{Question 7.}

\begin{parts}
  \item What is a pure virtual function and an abstract class? Write a program demonstrating the use of virtual destructors. \pts{7}
  \item What is exception handling? Explain the use of   \texttt{try},   \texttt{catch}, and   \texttt{throw} blocks in C++ with a sample program. \pts{7}
\end{parts}

\vfill
\rule{\linewidth}{0.4pt}

\begin{center}\small
\href{https://bhu-exam.vercel.app/}{Click here to visit our website}\\[4pt]\href{https://me-aryan.vercel.app/}{Click here to visit our contributor}\\[4pt]\href{https://quashan-me.vercel.app/}{Click here to visit our contributor}
\end{center}
\end{document}